\section{YOLO object detection neural network}

General purpose object detection is a much harder problem than image classification. In addition to the object class also the object location has to be found and in general there are more than one object in an image. Currently object detection systems are based on object classifiers applied to regions in the image. The simplest and computationally most expensive way is the sliding window approach. Here the object recognition is done on sub-images of different size at evenly spaced locations. More sophisticated approaches, like R-CNN, generate region proposals in a first step. On those regions the object recognition is done. Afterwards post-processing has to be done, like refinement of the bounding boxes and elimination of duplicate detections. Such R-CNNs are much faster as the sliding window implementations, especially with further optimizations like sharing features in different stages [https://papers.nips.cc/paper/5638-faster-r-cnn-towards-real-time-object-detection-with-region-proposal-networks.pdf]. However these systems are still not fast enough for real-time processing, at least on non-high-end hardware.

YOLO (You Only Look Once) puts all steps together into one neural network, therefore its name. The most recent version of YOLO is YOLOv2, which is an improved version of the original YOLO network and from here one YOLO refers to the YOLOv2-version. 

The image is segmented into a 13x13 grid. A grid cell is responsible for an object if the center of the object falls into the grid cell. Each grid cell predicts 5 bounding boxes, making a total of 845 bounding boxes. For each bounding box five values are predicted, x, y, w, h and confidence. The position of the bounding box is represented by the x- and y-coordinates of the box center. The size of the bounding box is given by the width w and the height h. The confidence is a value between 0 and 1 representing the intersection over union (IOU) of the bounding box and the ground truth. 

In addition to the bounding boxes the probability distribution over all object classes is calculated for every cell. The number of different object classes depends on the used training dataset. In this thesis YOLO trained on the COCO dataset was used. This training data contains more than 100,000 images with labels of 80 different object classes. Therefore the final prediction of the used YOLO version is a 13x13x105 tensor. 

To get the final result of the object detection non-maximum suppression is applied to get rid of bounding boxes responsible of objects being already better represented by another bounding box. In a last step the confidence of each bounding box is multiplied with the classification results of the corresponding cell to receive $< 845$ bounding boxes with object probabilities for each object class. 

